\paragraph{2. Preliminaries: Notation.}
Throughout this paper, $[d]$ denotes the set $\{1,2,\ldots,d\}$. For a random vector (or bit-string) $z$ in $\mathbb R^n$, we will use $z_i$ or $z(i)$ to denote the $i$-th coordinate of $z$, depending on whether we need to use the subscript for another index. We shall use $\mathbf i=(i_1,\ldots,i_d)$ for a $d$-tuple of indices. For a $d$-tuple $\mathbf i$, we write $|\mathbf i|=d$ to denote the size of the tuple.
For a matrix $M\in\mathbb C^{N\times N}$, we denote by $M^*$ its conjugate transpose. Given a string $x\in\mathbb R^N$, the $N\times N$ diagonal matrix with $x$ on the diagonal is denoted by $\operatorname{Diag}(x)$. The normalized trace of an $N\times N$ matrix $M$ is defined as $\operatorname{tr}_N(M)=\frac1N(\sum_{i=1}^NM_{ii})$. The operator norm of a matrix $M$ is denoted by $\|M\|_{\mathrm{op}}$. The left (resp. right) polar decomposition $(V,P)$ of a square matrix $M$ is a factorization of the form $M=VP$ (resp. $M=PV$) where $V$ is a unitary matrix and $P$ is a positive semidefinite matrix --- such a factorization always exists for any square matrix, and can be obtained easily from the singular-value decomposition of $M$. An $N\times N$ matrix $U$ is called a Haar random unitary if it is distributed according to the Haar measure on the Unitary group $U(N)$.
Random variables are typically denoted by capital letters (e.g., $X$). We write $\mathbb E[f(X)]$ and $\operatorname{Var}[f(X)]$ to denote the expectation and variance of the random variable $f(X)$; if $f:\{\pm1\}^m\to\mathbb R$ then we abbreviate it to $\mathbb Ef$ and $\operatorname{Var}[f]$, where the expectation and variance are taken with respect to the uniform measure on the discrete cube $\{\pm1\}^m$.
\paragraph{Fourier Analysis on the Discrete Cube.}
We give some basic facts about Fourier analysis on the discrete cube and refer to the book [28] for more details. Every function $f:\{\pm1\}^m\to\mathbb R$ can be written uniquely as a sum of monomials $\chi_S(x)=\prod_{i\in S}x_i$,
\[f(x)=\sum_{S\subseteq[m]}\widehat f(S)\chi_S(x),\tag{5}\]
where $\widehat f(S)=\mathbb E[f(X)\chi_S(X)]$ is the Fourier coefficient with respect to the uniform $X\in\{\pm1\}^m$. The monomials $\chi_S(x)=\prod_{i\in S}x_i$ form an orthonormal basis for real-valued functions on $\{\pm1\}^m$, called the Fourier basis. Parseval's identity implies that for uniform $X\in\{\pm1\}^m$,
\[\mathbb Ef^2=\sum_{S\subseteq[m]}\widehat f(S)^2,\qquad\operatorname{Var}[f]=\sum_{S\subseteq[m]:S\neq\varnothing}\widehat f(S)^2.\]
For a function on the hypercube, we define $\|f\|_2^2:=\mathbb Ef^2$ which can also be viewed as the sum of squared Fourier coefficients because of Parseval's identity.
The discrete derivative of a function on the hypercube $\{\pm1\}^m$ is given by
\[\partial_i f(x)=\frac12(f(x^{i\to1})-f(x^{i\to-1})),\]
where $x^{i\to b}$ is the same as $x$ except that the $i$-th coordinate is set to $b$. It is easily checked that $\partial_i f(x)$ coincides with the real partial derivative $\partial/\partial x_i$ of the real multilinear polynomial given by (5).
For a real-valued function $f:\{\pm1\}^m\to\mathbb R$, the influence of a variable $x_i$ on $f$ is defined as
\[\operatorname{Inf}_i(f)=\mathbb E|\partial_i f|^2=\sum_{S\subseteq[m]:i\in S}\widehat f(S)^2.\]
\paragraph{Block-multilinear Forms.}
A degree-$d$ block-multilinear form $f:\{\pm1\}^{n\times d}\to\mathbb R$ is given by
\[f(\mathbf x_1,\ldots,\mathbf x_d)=\mathbb Ef+\sum_{m=1}^d\sum_{\substack{|\mathbf b|=m\\|\mathbf i|=m}}\widehat f_{\mathbf b,\mathbf i}\,x_{b_1}(i_1)x_{b_2}(i_2)\cdots x_{b_m}(i_m),\tag{6}\]
where the tuple $\mathbf b=(b_1,\ldots,b_m)$ satisfies $1\leq b_1<\cdots<b_m\leq d$ and $\mathbf i\in[n]^m$ is an $m$-tuple. The expectation that yields the constant term is uniform over $\{\pm1\}^{n\times d}$. Note that $m$ is determined from the size of the tuple $\mathbf b$, so we just write $\widehat f_{\mathbf b,\mathbf i}$ above.
From Parseval's identity, the variance of $f$ and the influence of $x_b(i)$ on $f$ (where $b\in[d]$ and $i\in[n]$) are respectively given by
\[\operatorname{Var}[f]=\sum_{m=1}^d\sum_{\substack{|\mathbf b|=m\\|\mathbf i|=m}}\widehat f_{\mathbf b,\mathbf i}^{\,2},\qquad\operatorname{Inf}_{b,i}(f)=\sum_{m=1}^d\sum_{\substack{|\mathbf b|=m:b\in\mathbf b\\|\mathbf i|=m:i_b=i}}\widehat f_{\mathbf b,\mathbf i}^{\,2}.\]
From the above, it follows that for any block $b\in[d]$,
\[\sum_{i\in[n]}\operatorname{Inf}_{b,i}(f)^2\leq\operatorname{Var}[f]\leq\sum_{b\in[d],i\in[n]}\operatorname{Inf}_{b,i}(f)^2,\tag{7}\]
where the first inequality is an equality if $f$ is a homogeneous degree-$d$ block-multilinear form. For any $b\in[d]$, we write $\operatorname{MaxInf}_b(T)=\max\{\operatorname{Inf}_{b,i}(T):i\in[n]\}$ to denote the maximum influence of any variable in the block $\mathbf x_b$.
Since each non-constant monomial contains at most one variable from each block and the ordering of the blocks is fixed, a degree-$d$ block-multilinear form $f:\{\pm1\}^{n\times d}\to\mathbb R$ can be naturally viewed as a non-commutative polynomial in matrix variables with the constant term replaced with $\mathbb Ef$ times the identity. Denoting the non-commutative polynomial as $f(\mathbf U_1,\ldots,\mathbf U_d)$ where each $\mathbf U_b=(U_b(1),\ldots,U_b(n))$ is a block of non-commutative variables, the completely bounded norm $\|f\|_{\mathrm{cb}}$ of the form $f$ is defined as
\begin{multline*}
\|f\|_{\mathrm{cb}}=\sup\{\|f(\mathbf U_1,\ldots,\mathbf U_d)\|_{\mathrm{op}}:N\in\mathbb N,\ U_b(i)\in\mathbb C^{N\times N},\\
\|U_b(i)\|_{\mathrm{op}}\leq1,\ b\in[d],\ i\in[n]\}.
\end{multline*}
The completely bounded norm originates in the theory of operator algebras. In the literature, this norm is sometimes defined for homogeneous block-multilinear forms only, but here we extend the definition to non-homogeneous block-multilinear forms.
The supremum above is always attained and can be computed by solving a semidefinite program as shown by Gribling and Laurent [20]. One can also equivalently restrict the supremum in the definition above to unitary matrices since the convex hull of the set of unitary matrices is the unit operator-norm ball. Moreover, $\|T\|_\infty\leq\|T\|_{\mathrm{cb}}$ where $\|T\|_\infty=\max_{x\in\{\pm1\}^{n\times d}}|T(x)|$, so forms that are completely bounded are also bounded on the hypercube.
Note that if $f$ is a degree-$d$ block-multilinear form and if we fix some of the input bits to $\pm1$, then the resulting function $g$ is also a degree-$d$ block-multilinear form with the same blocks $\mathbf x_1,\ldots,\mathbf x_d$, but it does not depend on the variables that were fixed. It is also easy to see that $\|g\|_{\mathrm{cb}}\leq\|f\|_{\mathrm{cb}}$ because while computing $\|g\|_{\mathrm{cb}}$ we may restrict the matrix variables to $\pm I$ if that particular variable was set to $\pm1$. In other words, completely bounded norm does not increase under restrictions.
We write $\operatorname{MaxInf}(f)=\max_{b\in[d],i\in[n]}\operatorname{Inf}_{b,i}(f)$ for the maximum influence of any variable in $f$.
